\chapter*{Glossário}
\addcontentsline{toc}{chapter}{Glossário}
\label{chap:glossario}

\small

\noindent Este glossário reúne os termos técnicos e acrônimos utilizados ao longo do guia.

\bigskip

\noindent\textbf{ANDE} — \textit{Administración Nacional de Electricidad}. Empresa pública responsável pela distribuição de energia elétrica no Paraguai.

\medskip
\noindent\textbf{Aquífero Guarani} — Sistema aquífero transfronteiriço que se estende pelo Brasil, Argentina, Uruguai e Paraguai, com volume estimado em 37.000 km³ de água doce. Uma das maiores reservas de água subterrânea do mundo. No Paraguai, cobre a maior parte da região oriental.

\medskip
\noindent\textbf{Bortle (Escala)} — Escala de 1 a 9 que mede o brilho do céu noturno causado pela poluição luminosa. Bortle 1 = céu totalmente escuro (rural remoto); Bortle 9 = centro urbano com céu completamente iluminado. Relevante para agricultura (ritmos circadianos de animais e plantas) e qualidade de vida rural.

\medskip
\noindent\textbf{CONATEL} — \textit{Comisión Nacional de Telecomunicaciones}. Órgão regulador de telecomunicações do Paraguai, responsável pelo licenciamento de operadoras de celular e internet.

\medskip
\noindent\textbf{Censo 2022} — Censo Nacional de Población y Viviendas 2022, conduzido pela DGEEC. Primeira contagem completa desde 2002. Base para todos os dados demográficos deste guia.

\medskip
\noindent\textbf{DGEEC} — \textit{Dirección General de Estadística, Encuestas y Censos}. Instituto de estatística oficial do Paraguai, equivalente ao IBGE brasileiro.

\medskip
\noindent\textbf{Departamento} — Unidade administrativa de primeiro nível no Paraguai. São 17 departamentos mais o Distrito Capital (Assunção). Equivale ao estado brasileiro.

\medskip
\noindent\textbf{Distrito} — Unidade administrativa de segundo nível, equivalente ao município brasileiro. O Paraguai possui 262 distritos.

\medskip
\noindent\textbf{GSS} — \textit{Global Safety Score}. Pontuação de 0 a 10 desenvolvida para este guia, baseada na metodologia Skousen adaptada. Avalia a segurança e viabilidade de relocação de cada localidade em cinco dimensões ponderadas.

\medskip
\noindent\textbf{IDH} — Índice de Desenvolvimento Humano. Indicador composto (renda, saúde e educação) calculado pelo PNUD. No Paraguai varia de 0.648 (Caazapá) a 0.835 (Assunção).

\medskip
\noindent\textbf{INDERT} — \textit{Instituto Nacional de Desarrollo Rural y de la Tierra}. Órgão responsável pela reforma agrária e regularização fundiária rural. Principal fonte de dados sobre preço de terra rural.

\medskip
\noindent\textbf{ITAIPU} — Usina hidrelétrica binacional Brasil-Paraguai, com capacidade instalada de 14.000 MW. Supre praticamente toda a demanda elétrica do Paraguai e exporta o excedente para o Brasil.

\medskip
\noindent\textbf{Maquila} — Regime de manufatura no qual empresas estrangeiras instalam operações no Paraguai usando mão de obra local, com benefícios fiscais específicos (Lei 1.064/97). Responsável por uma parcela crescente das exportações paraguaias.

\medskip
\noindent\textbf{NASA POWER} — \textit{Prediction of Worldwide Energy Resource}. Base de dados climáticos da NASA derivada de modelos satelitais. Fornece irradiação solar, precipitação e outros parâmetros com cobertura global e série histórica de 20+ anos (2001--2020).

\medskip
\noindent\textbf{PETROPAR} — \textit{Petróleos Paraguayos}. Empresa estatal distribuidora de combustíveis. Seus postos de referência definem o preço-base da gasolina e diesel no país.

\medskip
\noindent\textbf{SAY} — \textit{Sistema Acuífero Yrenda-Toba-Tarijeño}. Aquífero da região do Chaco paraguaio. Apresenta água com alta concentração de sais, não adequada para consumo humano sem tratamento. Diferente do Aquífero Guarani da região oriental.

\medskip
\noindent\textbf{Skousen, Joel} — Autor norte-americano de \textit{Strategic Relocation: North American Guide to Safe Places} (1996, atualizado em 2011). Sua metodologia de avaliação de segurança por localidade foi adaptada para este guia como base do GSS.

\medskip
\noindent\textbf{SoilGrids} — Sistema global de mapeamento de solo desenvolvido pelo ISRIC (Centro Mundial de Informação do Solo). Fornece dados de pH, carbono orgânico, argila, areia e densidade do solo em grades de 250m de resolução.

\medskip
\noindent\textbf{Starlink} — Serviço de internet via satélite de baixa órbita da SpaceX. Disponível no Paraguai desde 2022, com planos residenciais a partir de USD 44/mês. Representa a principal alternativa de conectividade em regiões rurais remotas sem fibra ou rádio.

\medskip
\noindent\textbf{Terra Roxa} — Solo derivado do basalto da Formação Serra Geral, de coloração avermelhada, alta fertilidade natural e excelente estrutura física. Predominante no leste paraguaio (Alto Paraná, Itapúa, Canindeyú). Considerado um dos melhores solos agrícolas da América do Sul.

\medskip
\noindent\textbf{Tigo} — Maior operadora de telecomunicações do Paraguai (Millicom), com melhor cobertura 4G em áreas rurais. Recomendada como operadora principal para localidades fora dos centros urbanos.

\medskip
\noindent\textbf{Yacyretá} — Usina hidrelétrica binacional Argentina-Paraguai, no rio Paraná. Capacidade instalada de 3.200 MW. Complementa Itaipu no suprimento elétrico paraguaio.
